% !TEX root = ../main.tex
% !TEX program = lualatex
\documentclass[../main.tex]{subfiles}
\IfSubfilesClassLoaded{\externaldocument{\subfix{../../build/main}}}{}
% =============================================================================
% File: 22_AIT_Fundamentals_and_Lessons_Learned.tex
% =============================================================================
\begin{document}
\IfSubfilesClassLoaded{\chapter{EDO Study Guide}}{}
\section{AIT Fundamentals and Lessons Learned}

\subsection{Purpose and Advantages}
\begin{itemize}
  \item \textbf{Concentrated expertise.} An \ac{ait} is a trained unit (government, contractor, or hybrid) directed by an \ac{ait} manager to install specific modernization changes that require unique skills, tooling, or repetitive production-line execution~\autocite{TS9090-310,edo-4-1-4-ait-fundamentals-2025}.
  \item \textbf{Speed and repeatability.} Using the same team across hulls captures lessons learned, reduces quality escapes, and accelerates delivery of high-priority combat system upgrades (e.g., \ac{canes}, \ac{aws}, \ac{gpnts}) without overloading the shipyard workforce~\autocite{SL720-AA-MAN-030}.
  \item \textbf{Safety and configuration control.} Dedicated teams enforce technical work documents, boundary isolations, and logistics support for each alteration, reducing the risk of unauthorized changes that could violate \ac{subsafe} or Level I controls~\autocite{TS9090-310}.
\end{itemize}

\subsection{Roles and Responsibilities}
\begin{description}
  \item[\textbf{Sponsor / \ac{parm}.}] Assigns and funds the \ac{ait}, prioritizes hulls, and ensures logistics (allowance updates, training, tech manuals) are resourced before work begins~\autocite{SL720-AA-MAN-030}.
  \item[\textbf{\ac{ait} manager.}] Leads planning, safety, schedule integration, and quality execution; serves as the single point of contact for the Naval Supervising Activity and Ship's Force~\autocite{TS9090-310}.
  \item[\textbf{\ac{osic}.}] Coordinates day-to-day shipboard actions, manages tag-outs and system turnovers, and resolves emergent issues between the \ac{ait}, \ac{nsa}, and Ship's Force~\autocite{edo-4-1-4-ait-fundamentals-2025}.
  \item[\textbf{\ac{rmmco}.}] Screens alteration packages, aligns them with availability windows, and enforces waterfront integration and reporting requirements~\autocite{TS9090-310}.
  \item[\textbf{\ac{nsa}/\ac{lma}.}] Controls work authorization, boundary isolations, tests, and Objective Quality Evidence for work accomplished under the availability contract~\autocite{COMFLTFORCOMINST4790-3}.
\end{description}

\begin{figure}[htbp]
  \centering
  \begin{tikzpicture}[
    x=2.9cm,
    y=1.1cm,
    >=Latex,
    lane/.style={draw=Gray60, line width=0.6pt},
    swimstep/.style={draw=AccentBlue, fill=AccentBlue!10, rounded corners=2pt, align=center, inner sep=4pt, minimum width=2.6cm},
    arrow/.style={-Latex, line width=0.9pt, draw=Gray60}
  ]
    \foreach \y/\name in {0/Sponsor / PARM,-1/AIT Manager,-2/RMMCO,-3/NSA / LMA,-4/Ship's Force}{
      \draw[lane] (-0.4,\y) -- (4.3,\y);
      \node[left=3pt, align=right] at (-0.4,\y) {\small \name};
    }

    \node[swimstep] (need) at (0.6,0) {Screen need\\funding vector};
    \node[swimstep] (rmmco1) at (0.6,-2) {Package screen\\schedule slot};
    \node[swimstep] (plan) at (2.1,-1) {Plan install\\drawings / logistics};
    \node[swimstep] (moa) at (2.1,0) {Sign MOA\\funding + scope};
    \node[swimstep] (waf) at (2.1,-3) {Authorize work\\boundaries / WAF};
    \node[swimstep] (install) at (3.6,-1) {Execute install\\lead tests};
    \node[swimstep] (nsa) at (3.6,-3) {Witness tests\\accept OQE};
    \node[swimstep] (ship) at (3.6,-4) {Operate system\\sign turnover};

    \draw[arrow] (need) -- (rmmco1);
    \draw[arrow] (need) -- (plan);
    \draw[arrow] (rmmco1) -- (plan);
    \draw[arrow] (plan) -- (moa);
    \draw[arrow] (plan) -- (waf);
    \draw[arrow] (moa) -- (install);
    \draw[arrow] (waf) -- (install);
    \draw[arrow] (install) -- (nsa);
    \draw[arrow] (install) -- (ship);
    \draw[arrow] (nsa) -- (ship);
  \end{tikzpicture}
  \caption[AIT swim lane]{AIT swim-lane showing sponsor, AIT manager, RMMCO, NSA/LMA, and Ship's Force interactions from screening through certification.}
  \label{fig:ait-swimlane}
\end{figure}

\subsection{Planning Controls and MOA Discipline}
\begin{enumerate}
  \item \textbf{\ac{moa}.} The \ac{moa} outlines scope, funding, schedule, safety, habitability, quality, testing, and turnover; it must be signed by the \ac{ait} sponsor, \ac{rmmco}, Ship's Commanding Officer, and \ac{nsa} before work starts~\autocite{TS9090-310}.
  \item \textbf{Installation planning.} Each alteration requires a ship check, logistics certification, and completion of prerequisite changes; the \ac{nmp} databases confirm that planned work matches the approved Change Document~\autocite{SL720-AA-MAN-030}.
  \item \textbf{\ac{waf}/tag-out alignment.} Integration with the Work Authorization Form process and \ac{tum} prevents conflicting isolations between the \ac{ait} and shipyard trades~\autocite{S0400-AD-URM-010}.
  \item \textbf{Testing.} Entry/exit criteria, including pre-installation checkout, operational tests, and system certifications (e.g., Combat System Ship Qualification Test), are tailored per alteration and witnessed by \ac{nsa}~\autocite{TS9090-310}.
\end{enumerate}

\subsection{Lessons Learned for Board Prep}
\begin{itemize}
  \item \textbf{Late drawing changes drive rework.} Lock engineers early—\acp{ait} report that late drawing revisions, missing Interface Control Documents, or mismatched configuration data cause most delays~\autocite{edo-4-1-4-ait-fundamentals-2025}.
  \item \textbf{Habitability and safety must be priced.} Many \acp{ait} arrive without berthing or industrial hygiene plans; ensuring the \ac{moa} captures services, hot work permits, and fall protection prevents emergent cost growth~\autocite{SL720-AA-MAN-030}.
  \item \textbf{Document close-out.} Uploading \ac{oqe}, system configuration updates, and post-install logistics changes to the \ac{nde} is mandatory—omitting them erodes Fleet supply support~\autocite{TS9090-310}.
\end{itemize}

\IfSubfilesClassLoaded{
  \RenewDocumentCommand{\entryname}{}{\textbf{\color{Modern} Acronym}}
  \RenewDocumentCommand{\descriptionname}{}{\textbf{\color{Modern} Definition}}
  \printnoidxglossary[
    type=\acronymtype,
    title=Acronyms,
    style=long-booktabs
  ]}{}
\end{document}
